% Internal working draft; no venue selected.
% Synthetic teaching material, not an empirical or novelty claim.
% Passage requires amsmath and amssymb in the enclosing document.
\subsection{Updating a linear predictor while controlling changes on fixed examples}
\label{sec:prediction-retention}

How can a linear predictor incorporate a candidate update while retaining
its predictions on a fixed set of old examples? Let
\(X\in\mathbb{R}^{n\times d}\) contain the old feature vectors as rows,
and let \(w\in\mathbb{R}^d\) be the current parameter vector. The old
predictions are \(Xw\). We are given a candidate update
\(u\in\mathbb{R}^d\); the learning step that produces it is outside this
passage. The key observation is that a direction in the null space
\(\ker(X)=\{z\in\mathbb{R}^d:Xz=0\}\) cannot change these predictions.
We therefore project the candidate update onto that space before adding
it to the parameters.

\paragraph{Exact retention.}
For any \(X\in\mathbb{R}^{n\times d}\) and
\(w,u\in\mathbb{R}^d\), let \(r=\dim\ker(X)\) and let
\(N\in\mathbb{R}^{d\times r}\) have orthonormal columns spanning
\(\ker(X)\). In exact arithmetic, the update
\[
    v=NN^{\mathsf T}u,\qquad w'=w+v
\]
satisfies \(Xw'=Xw\).
Indeed, \(XN=0\), so
\(Xw'=Xw+XN N^{\mathsf T}u=Xw\).
If \(r=0\), \(N\) has no columns and \(NN^{\mathsf T}=0\), giving
\(v=0\). This elementary identity explains the construction: only
components of the update invisible to the fixed old examples survive.
Projection can also remove a direction needed to learn the new task, so
retention does not guarantee new-task improvement.

\paragraph{Thresholded implementation.}
The implementation uses a tolerance \(\epsilon>0\) and an orthogonal
projector \(P_\epsilon\in\mathbb{R}^{d\times d}\) onto the
right-singular subspace of \(X\) associated with singular values at most
\(\epsilon\), including the entire exact null space. It applies
\(v_\epsilon=P_\epsilon u\). This admits directions to which the old
predictions may have small, nonzero sensitivity. For every such
\(X,w,u,\epsilon\), the singular value decomposition gives, in exact
arithmetic,
\[
    \|XP_\epsilon\|_2\leq\epsilon,
    \qquad
    \|X(w+v_\epsilon)-Xw\|_2
    =\|XP_\epsilon u\|_2
    \leq\epsilon\|u\|_2.
\]
Here \(\|\cdot\|_2\) denotes the Euclidean norm for vectors and the
spectral operator norm for matrices. Thus the implementation bounds the
change in the old prediction vector rather than necessarily preserving
it exactly; the bound scales with the candidate update's magnitude.

For an exact-arithmetic illustration, take
\(X=\operatorname{diag}(1,0.01)\),
\(w=(2,3)^{\mathsf T}\), \(u=(1,4)^{\mathsf T}\), and
\(\epsilon=0.02\). The exact null space is trivial, so the idealized
update is zero. The thresholded projector retains the second coordinate,
yielding \(v_\epsilon=(0,4)^{\mathsf T}\) and prediction shift
\((0,0.04)^{\mathsf T}\).

Both statements concern predictions on the rows of this fixed \(X\).
They do not establish task accuracy, retention on other examples, or
generalization. The illustration is not empirical evidence, and actual
floating-point numerical error has not been separately bounded.
